\section{Related Work}
\label{sec:related}

\textbf{Vision-language-action models.} Generalist policies increasingly couple a pretrained vision-language backbone to an action decoder trained on large cross-embodiment datasets: \mbox{RT-2}~\cite{brohan2023rt2} grounds web knowledge in control, Octo~\cite{octo2024} and OpenVLA~\cite{kim2024openvla} are open generalist policies, and $\pi_0$~\cite{black2024pi0} adds a flow-matching action expert. These models are large and data-hungry; SmolVLA~\cite{shukor2025smolvla}, used here, keeps the flow-matching design but targets a single consumer GPU and community-scale data. Such single-skill competence does not, however, guarantee reliable long-horizon execution on a fixed low-cost arm --- the gap this work addresses.

\textbf{Imitation learning and its failure modes.} These policies are trained by behavioural cloning, as are ACT~\cite{zhao2023act} and diffusion policies~\cite{chi2023diffusion}. Cloning from successful demonstrations is prone to drift: small errors carry the policy into states absent from training, where mistakes compound --- the covariate-shift problem behind interactive correction like DAgger~\cite{ross2011dagger}. How this error scales with horizon is still debated~\cite{spencer2021feedback,foster2024horizon}, but in this setting the consequence is concrete: a single fumbled grasp can leave the arm in a configuration the policy never recovers from. This motivates adding structure \emph{around} the policy rather than enlarging it.

\textbf{Decomposition and language hierarchies.} A common remedy breaks a long task into language-specified pieces. RT-H~\cite{belkhale2024rth} inserts a language-motion layer between task and action; SayCan~\cite{ahn2022saycan} grounds an instruction in a library of affordance-scored skills; and planners such as Inner Monologue~\cite{huang2022inner} and Code-as-Policies~\cite{liang2022cap} sequence or generate skills from language. These delegate high-level control to a (often large) language model and typically assume each invoked skill succeeds. The orchestrator introduced here keeps the language-conditioned subtask decomposition but replaces the language-model planner with a lightweight finite-state controller that verifies each subtask's physical outcome from privileged state, making intermediate failure observable rather than assumed away.

\textbf{Failure detection and recovery.} Recent work reacts to execution failures: RACER~\cite{dai2024racer} augments demonstrations with recovery trajectories supervised by a VLM; SAFE~\cite{gu2025safe} learns a task-agnostic failure detector for VLAs; Code-as-Monitor~\cite{zhou2024codeasmonitor} and Sentinel~\cite{agia2024sentinel} monitor execution at runtime; and REFLECT~\cite{liu2023reflect} explains failures with a language model for replanning. These either need failure-labelled data or add a learned/language-model monitor atop the policy. This work instead defines recovery \emph{structurally}, at two levels of the goal hierarchy --- retrying the same target (Level-1) versus redirecting to a different feasible target (Level-2) --- adjudicated by privileged physical checks.

\textbf{Self-improvement from autonomous data.} Finally, the autonomous data loop relates to policy self-improvement: keeping only successful self-collected trajectories and retraining is reward-filtered self-imitation --- the binary-reward limit of reward-weighted regression~\cite{peters2007rwr} --- akin to autonomous-practice systems like MEDAL++~\cite{sharma2023selfimproving}. The orchestrator's privileged checks supply the filter; \cref{sec:experiments} reports when this helps and when it does not.
